% ===========================================================================
\section{Discussion}
\label{sec:discussion}
% ===========================================================================

\subsection{The methane correction factor, and why it is not resolved here}

The sewer methane correction factor is the largest single assumption in this sector and the
one the evidence does not settle. It deserves the space because how it is handled is a
methodological position, not a parameter choice.

IPCC gives \num{0.5} for stagnant open sewers and \num{0.1} for wastewater discharged to a
flowing, well-aerated body. The distinction was drawn for systems that are clearly one or the
other. Pakistani urban drainage is not: the drains are open, warm, slow-moving and heavily
organically loaded, which is the description \num{0.5} exists for, while a substantial
fraction ultimately discharges to rivers and irrigation channels, which is the description
\num{0.1} exists for. Both descriptions are accurate about different parts of the same system,
and no Pakistani source reports the hydraulic condition of the sewered fraction, which is the
quantity the factor is actually asking about.

The consequence is a factor of \num{1.59} on the sector's largest category:
\SI{24.46}{\mega\tonne} against \SI{15.39}{\mega\tonne}, or \SI{37.45}{\mega\tonne} against
\SI{28.38}{\mega\tonne} for the sector.

\subsubsection{Why external agreement is not a tiebreak}

At \num{0.1} this sector lands close to EDGAR, and it would be easy to present that as
corroboration. It is not, and treating it as one would be a methodological error worth naming
because it is a tempting one.

EDGAR is a model. Its Pakistani wastewater figure rests on its own assumption about the same
unmeasured quantity---the anaerobic fraction of Pakistani sewage---made from the same absence
of Pakistani measurement. If this inventory adopted \num{0.1} because it produces agreement,
the two would agree because both had chosen the same default, and the agreement would then be
cited as though it were evidence about Pakistani drains. That is how a shared assumption
becomes an apparent consensus.

We therefore report \num{0.5} as central, because the physical description of the drains is
the stronger evidence and is specific to Pakistan, and publish \num{0.1} beside it so a user
who weighs the comparator differently can take it. What we do not do is let the comparison
choose.

\subsubsection{What would settle it}

Dissolved-methane concentration and flux measurement on urban drains in two or three Pakistani
cities, sampled across a seasonal cycle to capture the dry-season stagnation and the monsoon
flushing that the two IPCC categories describe separately.

The case for doing it is quantitative rather than rhetorical. Domestic wastewater is
\SI{5.5}{\percent} of Pakistan's national emissions and \SI{13.0}{\percent} of national
inventory variance under Approach~1---the second largest variance contributor of any category
in the country on that screen.
A measurement campaign of a few tens of thousands of dollars would reduce national inventory
uncertainty more than any comparable expenditure elsewhere in the waste or energy sectors.

\subsection{What the closed mass balance buys, and what it does not}

The balance guarantees that the collected and uncollected streams together account for
exactly the waste the country generates, with the burned quantity a declared fraction of the
uncollected stream and the unburned \SI{6.27}{\mega\tonne} carried as an explicit unresolved
fate rather than assigned to either category. That removes a specific class of error---the two
categories drifting into mutual inconsistency as one is updated---and makes the collection rate
an explicit, single, arguable number rather than an implicit consequence of two independent
estimates.

It is worth being precise about what it does not buy, because a closed balance invites
overconfidence. It says nothing about whether the split point is right. The collection rate of
\SI{50.3}{\percent} carries real uncertainty, and moving it moves waste between two fates with
quite different methane yields per tonne: a tonne landfilled decays slowly and anaerobically
over decades, while a tonne burned emits immediately and mostly as \COtwo{}. Nor does it
validate the generation figure itself; a partition of a wrong total still closes.

What it changes is the shape of the argument. A reader who disagrees with the collection rate
can see exactly which two categories move and in which directions, rather than having to work
out whether two independently sourced numbers are mutually consistent.

\subsection{Implications}

Three follow, in descending order of confidence.

\textbf{Wastewater is the sector, and the sector is a measurement problem.} Two-thirds of the
emissions and nearly all of the tractable uncertainty sit in one category whose central
parameter has never been measured in Pakistan. Effort spent refining the solid-waste decay
parameters or the burning composition would improve smaller numbers that are already better
constrained.

\textbf{Open burning matters for air quality, not for climate.} At \SI{0.86}{\percent} of
national \COe{} and roughly \SI{14}{\percent} of national primary \PM{}, it is a category whose case for
intervention rests on exposure. An abatement programme justified on greenhouse-gas grounds
would be hard to defend; the same programme justified on particulate would be
straightforward.

\textbf{The collection rate is a policy variable, not just a parameter.} Because the partition
makes it the single number governing what fraction of waste burns rather than decays,
improving collection is visible in this inventory as a movement between categories. That is
the right structure for evaluating a municipal collection programme, and it is not available
in an inventory whose two categories are estimated independently.
